\documentclass[12pt,a4paper]{article}

\usepackage[utf8]{inputenc}
\usepackage[T1]{fontenc}
\usepackage[french]{babel}
\usepackage[a4paper,margin=2.2cm]{geometry}
\usepackage{enumitem}
\usepackage{tcolorbox}
\usepackage{booktabs}
\usepackage{xcolor}

\setlist[itemize]{nosep}
\setlist[enumerate]{nosep}

\title{\textbf{Correction : TD -- Étude de cas : Perte d’accès Internet}\\
\large BTS SIO -- Gestion d’incident \& Réseaux}
\author{Maël}
\date{\today}

\begin{document}

\maketitle

\section*{Scénario 1 -- Poste isolé sans réseau}
\begin{enumerate}
    \item \textbf{Cause probable :} Échec de récupération d'une configuration IP via le serveur DHCP. L'adresse en 169.254.x.x indique que le poste s'est auto-attribué une adresse.
    \item \textbf{Localisation :} Entre le poste client et le switxh d'accès ou le serveur DHCP (câble débranché, port défectueux).
    \item \textbf{Action corrective :} Vérifier le branchement du câble Ethernet, tester avec un autre câble, ou forcer une nouvelle d'IP.
    \item \textbf{Priorité :} Basse. Un seul utilisateur est impacté.
\end{enumerate}




\section*{Scénario 2 -- Problème de résolution DNS}
\begin{enumerate}
    \item \textbf{Réseau local :} Oui, le ping vers la passerelle est OK.
    \item \textbf{Internet accessible :} Oui, le ping vers l'IP 8.8.8.8 (Google) répond. La connectivité IP est établie.
    \item \textbf{Type de problème :} Problème de résolution de DNS. Le système ne parvient pas à traduire "google.fr" en adresse IP.
    \item \textbf{Solution :} Vérifier la configuration des serveurs DNS sur le poste ou le serveur DHCP. Tester avec un DNS public (1.1.1.1) pour confirmer.
\end{enumerate}




\section*{Scénario 3 -- Incident global}
\begin{enumerate}
    \item \textbf{Nature :} Incident global (plusieurs utilisateurs impactés).
    \item \textbf{Localisation :} Équipement d'interconnexion central (Routeur ou Switch cœur de réseau).
    \item \textbf{Priorité :} Critique. L'activité de l'entreprise est paralysée.
    \item \textbf{Escalade :} Oui, escalade fonctionnelle vers l'administrateur réseau (Niveau 2 ou 3) car le matériel semble hors service physiquement.
\end{enumerate}




\section*{Scénario 4 -- Panne fournisseur}
\begin{enumerate}
    \item \textbf{Réseau interne :} Oui, le ping passerelle et les services internes fonctionnent.
    \item \textbf{Localisation :} Extérieure à l'entreprise. Lien WAN ou équipement de l'Opérateur.
    \item \textbf{Action :} Contacter le support technique du fournisseur d'accès (Orange, SFR, etc.) pour ouvrir un ticket d'incident sur la ligne.
    \item \textbf{Communication :} Informer les utilisateurs qu'une coupure généralisée d'Internet est en cours suite à une panne chez l'opérateur et donner un délai estimé de rétablissement si possible.
\end{enumerate}
\section*{Synthèse finale (Scénario 3)}
\vspace{0.5cm}

\noindent Dans un premier temps, création d'un ticket d'incident majeur, prise d'informations sur l'incident. Vérification de la passerelle et des branchement. 
\vspace{0.3cm}

\noindent Une Communication à tous les utilisateurs de la panne. Et recherche de slution temporaire pour continuer l'activité (partage de connection temporaire, redirection vers un autre site, etc).
\vspace{0.3cm}

\noindent Renseignement des infos et actions réalisées dans le tickets pour le suivi (test effectués, communications aux utiisateurs, etc) puis escalade vers le niveau 2 ou 3 pour une intervention sur le matériel.
\vspace{0.3cm}

\noindent Enfin, une fois le probleme resolu, une communication de fin d'incident et un rapport d'incident détaillé pour l'analyse.
\vspace{0.3cm}

\noindent Pour finir, proposition de cloture du ticket (verifier que tout conviens aux utilisateurs) avec une mise en place de mesures correctives pour eviter le retour du probleme.

\end{document}